\section{Using the estimate for an actual firing}\label{sec:interpretation}
The nominal calculations give roughly 1.7--1.8 minutes of continuous
exposure for the 4.7625 mm sustainer liner and 2.3--2.5 minutes for the
6.350 mm booster liner. This clock starts at a particular wall station
when gas first reaches bare phenolic. An inhibitor, a covering grain or a
local breach would change that starting time. The available grain length and
stoichiometric reaction alone determine neither when that station opens nor
the heat flux it receives.

\paragraph{Establish the load at the liner.}
The tested gas-film coefficients span a factor of four. At fixed nominal
pressure, the CF booster changes from about 259 to 75 s across that span.
Record a time-resolved thermal measurement at a known axial station,
the chamber-pressure trace and when the liner there becomes exposed.
A measured heat-flux history can replace the assumed constant gas-side load.
A rear-face thermocouple history supplies a separate check on heat conducted
through the wall. Section~\ref{sec:env} shows the radiative load explicitly.

\paragraph{Measure what the phenolic loses.}
For the Al sustainer at the nominal film coefficient, halving the assumed
char-removal energy gives about 72 s and doubling it gives about 185 s.
Under a known heat load, measure the specimen mass, surviving liner thickness,
char depth and temperatures at recorded exposure times. A series of durations
separates conversion within the liner from loss of the exposed char.
Use those measurements to fit the removal law and reserve a separate exposure
for checking the fitted prediction.

\paragraph{Separate liner exhaustion from loss of pressure containment.}
The Al calculation holds melted metal in place, and the CF calculation
continues beyond the 1000~$^\circ$C limit of the adopted transport properties.
The temperature histories on page~\pageref{sec:temperatures} give the
thermal input for a separate assessment using the actual alloy, laminate,
bondline and their hot strength. A local crack or debond can expose the casing
while a uniform sector still predicts surviving liner.

\paragraph{Where to inspect the calculation.}
The \href{https://github.com/MDNich/MIT-RT-26_27-Work/tree/main/simulation_work/solidprop/iris/v1/output/data}{two nominal history files}
separate ANSYS samples from the independent finite-volume inventories.
The \href{https://github.com/MDNich/MIT-RT-26_27-Work/tree/main/simulation_work/solidprop/iris/v1/validation}{validation records}
give the last time intervals, mesh comparisons, mass and heat balances and input
hashes. The \href{https://github.com/MDNich/MIT-RT-26_27-Work/blob/main/simulation_work/solidprop/iris/v1/output/pdf/iris_phenolic_report.pdf}{companion report}
prints the historical MAPDL and Fortran material routines. Appendix~B here
prints the continuous-exposure MAPDL input and its thermal-step generator.
\clearpage
